\section{Algebra, infinity, and formalization}

\begin{frame}{Operations are constrained relations}
  \bookref{Chapter 4, Sections 1--3}
  A binary operation $\mu:X\times X\to X$ is encoded by
  \[
    \Gamma_\mu=\{(x,y,z)\in X^3:\mu(x,y)=z\}.
  \]
  The relation is functional and total in the third coordinate:
  \[
    \forall x,y\in X\;\exists!z\in X:\;(x,y,z)\in\Gamma_\mu.
  \]

  \begin{theorem}
    A permutation $f:X\to X$ preserves $\mu$ iff it preserves $\Gamma_\mu$:
    \[
      f(\mu(x,y))=\mu(f(x),f(y))
      \iff
      f^{\times3}(\Gamma_\mu)=\Gamma_\mu.
    \]
  \end{theorem}

  Thus finite groups, rings, semigroups, lattices, and general universal
  algebras enter the same automorphism framework.
\end{frame}

\begin{frame}{A complete two-element operation ensemble}
  \bookref{Chapter 4, Section 4}
  There are $n^{n^2}$ labeled binary operations on an $n$-element set; for
  $n=2$, this is $16$.

  \begin{columns}[T]
    \column{0.48\textwidth}
      \centering
      \begin{tabular}{ccc}
        \toprule
        $|\Aut|$ & classes & labeled tables \\
        \midrule
        $1$ & $6$ & $12$ \\
        $2$ & $4$ & $4$ \\
        \midrule
        total & $10$ & $16$ \\
        \bottomrule
      \end{tabular}
    \column{0.48\textwidth}
      \centering
      \begin{tabular}{lcc}
        \toprule
        constraint & labeled & classes \\
        \midrule
        associative & $8$ & $5$ \\
        commutative & $8$ & $4$ \\
        monoid & $4$ & $2$ \\
        group & $2$ & $1$ \\
        \bottomrule
      \end{tabular}
  \end{columns}

  \vspace{0.7em}
  Algebraic laws define sharply smaller subensembles, while orbit multiplicity
  continues to be $n!/|\Aut(A)|$.
\end{frame}

\begin{frame}{In infinity, the signature becomes decisive}
  \bookref{Chapter 5, Sections 1--2}
  \small
  \centering
  \begin{tabular}{lll}
    \toprule
    underlying set & specified structure & residual symmetry \\
    \midrule
    $\Z$ & additive group & $n\mapsto\pm n$ \\
    $\Z$ & unital ring & identity only \\
    $\Z$ & ordered unital ring & identity only \\
    $\R$ & $\mathbb Q$-vector space & many linear automorphisms \\
    $\R$ & field & identity only \\
    $\R^n$ & vector space & $\mathrm{GL}_n(\R)$ \\
    $\R^n$ & inner-product space & $O(n)$ \\
    \bottomrule
  \end{tabular}

  \vspace{0.4em}
  Infinite groups of the same cardinality may represent radically different
  symmetry. Consequently,
  \[
    \log\frac{|\Aut(X)|}{|\Aut(A)|}
  \]
  is generally not a meaningful numerical invariant.
\end{frame}

\begin{frame}{Candidate infinite symmetry deficits}
  \bookref{Chapter 5, Sections 3--4; Chapter 6}
  \centering
  \begin{tabular}{lll}
    \toprule
    candidate & natural setting & principal limitation \\
    \midrule
    $\log[G:H]$ & finite-index residual symmetry & often infinite \\
    $\dim G-\dim H$ & Lie/algebraic groups & misses wild groups \\
    quotient measure & compact/measured groups & needs invariant measure \\
    finite-window rate & growing physical domains & boundary dependent \\
    measure relative entropy & phase laws on $G/H$ & needs reference measure \\
    \bottomrule
  \end{tabular}

  \vspace{0.7em}
  \begin{block}{Reference-relative requirement}
    Choose a structured background $B$ and $G\leq\Aut(B)$. Compare the phase
    stabilizer $H$ with $G$; $G=\Sym(X)$ is reserved for a genuinely bare set.
  \end{block}
\end{frame}

\begin{frame}{Relational scaling taxonomy}
  \bookref{Chapter 6, Section 6}
  \small
  \begin{center}
    \begin{tabular}{lll}
      \toprule
      quantity & question answered & caution \\
      \midrule
      $I_G=\log[G:H]$ & which finite sector is selected? & may be subextensive \\
      $\widetilde I_G=I_G/\log|G|$ & what fraction of $G$ is lost? & not a density \\
      $I_L/a_L$ & is symmetry loss extensive? & declare $a_L$ \\
      $D(\rho),\widehat D$ & is a coarse pattern visible? & model-dependent proxy \\
      $a_L\operatorname{Var}(Q_L)$ & how do fluctuations scale? & depends on $Q_L$ \\
      \bottomrule
    \end{tabular}
  \end{center}

  \begin{block}{Reporting rule}
    Use total deficit for discrete selection and a declared deficit rate for
    extensive ordering. Never identify a block proxy with a normalized
    automorphism deficit without an explicit relation between them.
  \end{block}
\end{frame}

\begin{frame}{Machine-checked core}
  \bookref{Appendix D; Lean sources and formalization ledger}
  \footnotesize
  Every numbered statement is paired with a checked ledger entry.

  \begin{columns}[T,onlytextwidth]
  \begin{column}{0.43\textwidth}
    \scriptsize
    % Generated by scripts/check_lean_coverage.py --write; do not edit.
\begin{tabular}{lr}
\hline
numbered-statement status & count \\
\hline
complete & 46 \\
partial & 1 \\
needs definitions & 1 \\
analysis-heavy & 7 \\
external citation & 1 \\
compatibility condition & 1 \\
\hline
total & 57 \\
\hline
\end{tabular}

  \end{column}
  \begin{column}{0.54\textwidth}
  \begin{itemize}
    \item Chapters 1--2 and 4 connect automorphism subgroups, relabeling orbits,
      finite partition sums, and extremal equality cases.
    \item Chapter 3 checks general-mean $O(d^6)$ control, quartic local minima,
      the inverse-branch $3/8$ limit, and orthogonal field-energy invariance.
    \item Chapter 6 also checks Haar index, relative entropy, and scalar-unitary
      invariance of fixed-number and one-particle operators.
    \item Continuum limits, quotient geometry, gauge variation, and XY analysis
      remain open; the condensation theorem is an external citation.
  \end{itemize}
  \end{column}
  \end{columns}

  \medskip
  Appendix D records the proof scope. Two unnumbered identities and nine
  supplementary targets are tracked separately.
\end{frame}
